\begin{corollary}{5.5}\label{VIB.5.5}
Let $S$ be a scheme, $G$ an $S$-group scheme, locally of finite presentation, with connected fibers, and universally open over $S$. Then $G$ is separated and of finite presentation over $S$.\nde{Let us point out here the following result ([Ray70a], VI 2.5): if $S$ is normal, $G$ smooth with connected fibers, and $X$ is an (fppf) homogeneous $G$-space (i.e. the morphisms $X\to S$ and $G\times_SX\to X\times_SX$ are covering for the (fppf) topology), locally of finite type over $S$, then $X$ is locally quasi-projective over $S$. In particular, $G$ is quasi-projective over $S$. See also N.D.E. (35) in VIA.}
\end{corollary}
Indeed, by 3.6 and 5.4, $G$ is quasi-compact and separated over $S$, and, since $G$ is locally of finite presentation over $S$, it is therefore of finite presentation over $S$.

\numpara{5.6}\textbf{Proof of Theorem 5.3.} Before establishing 5.3, let us prove a few lemmas.
\begin{lemma}{5.6.0}\label{VIB.5.6.0}\nde{This lemma, communicated by O. Gabber, has been added; it improves EGA IV$_3$, 14.4.1.2 and corrects the proof of loc. cit., 14.4.1.3 without changing its hypotheses (compare with the erratum (ErrIV, 38) in EGA IV$_4$).}
(i) Let $A\subset B$ be integral domains, with $B$ integral over $A$, and let $\mathfrak q\in\Spec(B)$ be such that $A$ is unibranch at the point $\mathfrak p=\mathfrak q\cap A$. Then the morphism $\pi:\Spec(B)\to\Spec(A)$ is open at the point $\mathfrak q$.

(ii) Let $X,Y$ be two irreducible preschemes, $f:X\to Y$ a dominant morphism, $x$ a point of $X$ such that $f$ is quasi-finite at $x$ and $y=f(x)$ is a unibranch point of $Y$. Then $f$ is open at the point $x$. In particular, if $Y=\Spec(\OO_{Y,y})$ is a local prescheme with closed point $y$, then $f(U)=Y$ for every neighborhood $U$ of $x$.
\end{lemma}
(i) Let $K$ (resp. $L$) be the field of fractions of $A$ (resp. $B$), $A'$ the normalization of $A$, and $B'$ the subring of $L$ generated by $A'$ and $B$. Then $B'$ is integral over $A'$. Put $Y=\Spec(A)$, $X=\Spec(B)$, $Y'=\Spec(A')$, $X'=\Spec(B')$, so that one has a commutative diagram
\[
\xymatrix{X'\ar[r]^{\pi'}\ar[d]&Y'\ar[d]^\phi\\X\ar[r]^\pi&Y}
\]
in which all the morphisms are integral and surjective.

Since $A$ is unibranch at $\mathfrak p$, $Y'$ has a unique point $\mathfrak p'$ above $\mathfrak p$; consequently, if $U$ is an open neighborhood of $\mathfrak p'$ in $Y'$, $\phi(Y'-U)$ is a closed subset not containing $\mathfrak p$, so that its open complement is contained in $\phi(U)$. This shows that $\phi:Y'\to Y$ is open at $\mathfrak p'$, and it therefore suffices to show that $\pi'$ is open. One is thus reduced to the case where $A=A'$ is normal.

Let $N$ be a quasi-Galois extension of $K$ containing $L$, let $\widetilde X=\Spec(\widetilde B)$, where $\widetilde B$ is the integral closure of $B$ in $N$, and let $G=\Aut(N/K)$. Since $\widetilde X\to X$ is surjective, it suffices to show that $\widetilde\pi:\widetilde X\to Y$ is open. Let $U$ be an open subset of $\widetilde X$ and $U'=\bigcup_{g\in G}gU$. Since $G$ acts transitively on the fibers of $\widetilde X\to Y$ (cf. [BAC], Chap. V, §2.3, Prop. 6), $\widetilde\pi(U)$ equals $\widetilde\pi(U')$, and the latter equals the open complement of the closed subset $\widetilde\pi(\widetilde X-U')$. This proves (i).

(ii) One may suppose $Y$ and $X$ reduced. By ``Zariski's main theorem'' (cf. EGA IV$_3$, 8.12.9), there are affine open neighborhoods $U$ of $x$ and $V=\Spec(A)$ of $y$, such that $f(U)\subset V$, and a factorization:
\[
\xymatrix{U\ar@{^{(}->}[r]^j\ar[dr]_f&V'\ar[d]^u\\&V}\ ,
\]
where $j$ is an open immersion and $u$ is finite. Replacing $V'$ by the closure of $j(U)$, one may suppose $V'$ irreducible, hence $V'=\Spec(B)$, where $B$ is an integral $A$-algebra, finite over $A$. Moreover, since $f$ is dominant, so is $u$, so that the morphism $A\to B$ is injective. Since, by hypothesis, $A$ is unibranch at the point $y$, it follows from (i) that $u$ is open at the point $j(x)$, and hence $f=u\circ j$ is open at the point $x$. This proves the first assertion of (ii). The second follows from it, for, if $Y$ is a local prescheme with closed point $y$, every open subset containing $y$ equals $Y$. (In the case where $Y=\Spec(\OO_{Y,y})$, one may also use, instead of EGA IV$_3$, 8.12.9, the local form of Zariski's main theorem, found for example in [Pes66], or [Ray70b], Ch. IV, Th. 1.)
\begin{lemma}{5.6.1.0}\label{VIB.5.6.1.0}\nde{This lemma, used in the proof of Lemma 5.6.1, has been made explicit.}
Let $f:X\to S$ be a morphism locally of finite presentation, $s\in S$, and $x\in X_s$. Let $n=\dim_x(X_s)$ and let $q$ be the structure morphism $\mathbb A_S^n\to S$.

Suppose given a quasi-finite morphism $u:X\to\mathbb A_S^n$ such that $f=q\circ u$, and suppose $f$ universally open at the generic point $z$ of an irreducible component $\mathfrak z$ of $X_s$, containing $x$ and of dimension $n$. Then $u$ is universally open at the point $x\in X$.
\end{lemma}
First observe that $(\dagger)$ the hypotheses are preserved by every base change $\pi:S'\to S$ covering $s$ (i.e. such that $\pi^{-1}(s)\ne\varnothing$). Indeed, let $S'\to S$ be such a morphism, let
\[
\xymatrix{X'\ar[r]^{u'}\ar[dr]_{f'}&\mathbb A_{S'}^n\ar[d]^{q'}\\&S'}
\]
be the diagram obtained by base change, and let $s'$ be a point of $S'$ above $s$ and $x'$ a point of $X'_{s'}$ above $x$. Since $X'_{s'}=X_s\otimes_{\kappa(s)}\kappa(s')$, by lifting of generizations and invariance of dimension under field extension (EGA IV$_2$, 2.3.4 (i) and 4.1.4), $x'$ is contained in an irreducible component $\mathfrak z'$ of $X'_{s'}$ whose generic point $z'$ lies above $z$, and one has $n\le\dim\mathfrak z'\le\dim_{x'}X'_{s'}\le n$, whence $\dim\mathfrak z'=n=\dim_{x'}X'_{s'}$. Since $f$ is universally open at $x$, $f'$ is universally open at $x'$.
% typo?: kept as printed: x and x' in the last sentence rather than z and z'.

Put $Y=\mathbb A_S^n$. By EGA IV$_3$, 14.3.3.1 (i), to prove that $u$ is universally open at $x$, it suffices to prove that, for every integer $r\ge0$ and every point $x'$ of $X'=X[T_1,\ldots,T_r]$ above $x$, the morphism $u':X'\to Y'=\mathbb A_S^n[T_1,\ldots,T_r]$ is open at the point $x'$. Now, writing $S'=S[T_1,\ldots,T_r]$ and $q'$ for the projection $Y'\simeq\mathbb A_{S'}^n\to S'$, one is in the situation obtained by the base change $S'\to S$. Thus, by the preceding argument, one is reduced to showing that $u$ is open at the point $x$.

Put $y=u(x)$. Since $u$ is locally of finite presentation (since $f$ and $q$ are, cf. EGA IV$_1$, 1.4.3), by EGA IV$_1$, 1.10.3, it suffices to show that $u(\Spec\OO_{X,x})=\Spec(\OO_{Y,y})$. For this one may suppose $S$ affine and integral. Then let $\pi:S'\to S$ be its normalization; denote by $u':X'\to Y'$ the morphism obtained from $u$ by base change. Since the morphism $\pi_Y:Y'\to Y$ is integral and surjective, one has
\[
\Spec(\OO_{Y,y})=\bigcup_{y'}\pi_Y(\Spec(\OO_{Y',y'})),
\]
the union being taken over all the points of $Y'$ above $y$; it therefore suffices to show that, for each such $y'$, and every $x'\in X'_{y'}$ above $x$, one has $u'(\Spec\OO_{X',x'})=\Spec\OO_{Y',y'}$. Since the hypotheses are preserved by the base change $\pi:S'\to S$, one is thus reduced to the case of $S'$, i.e. one may suppose $S$ an integral normal scheme.
% typo: le le changement -> le changement.

Now the hypotheses on $f$ imply, by EGA IV$_3$, 14.3.13, that there is an irreducible component $Z$ of $X$ containing $z$ (and hence $x$), dominating $S$ and such that
\[
\dim_z(Z_s)=n=\dim(Z_\eta),
\]
where $\eta$ is the generic point of $S$. Let $\xi$ be the generic point of $Z$. Since $u$, hence also $u_\eta:Z_\eta\to\mathbb A_\eta^n$, is quasi-finite, the closure in $\mathbb A_\eta^n$ of the point $u_\eta(\xi)=u(\xi)$ has dimension $n$, hence $u(\xi)$ is the generic point of $\mathbb A_\eta^n$, which is also the generic point of $Y=\mathbb A_S^n$. Consequently, denoting by $g$ the restriction of $u$ to $Z$, the morphism $g:Z\to Y$ is quasi-finite and dominant. Since $Y$ is normal, it follows from Lemma 5.6.0 that $g$ is open, so $u$ is open at every point of $Z$, in particular at the point $x$. Consequently, one has $u(\Spec\OO_{X,x})=\Spec\OO_{Y,y}$, and this completes the proof of Lemma 5.6.1.0.

The following lemma advantageously replaces EGA IV$_3$, 14.5.10,\nde{The original, which indicated 19.5.10, has been corrected.} in that it is independent of noetherian hypotheses.
\begin{lemma}{5.6.1}\label{VIB.5.6.1}
Let $S$ be a scheme, $f:X\to S$ an $S$-scheme locally of finite presentation, $s$ a point of $S$, $x$ a closed point of $X_s$. Suppose $f$ universally open at the generic point $z$ of an irreducible component $\mathfrak z$ of $X_s$, containing $x$ and such that $\dim_x(X_s)=\dim\mathfrak z$. Then there is a commutative diagram:
\[
\xymatrix{X\ar[r]^f&S\\S''\ar[u]^h\ar[ur]^\phi\ar[r]^\pi&S'\ar[u]_w}\ ,
\]
where $S'$ is an affine scheme, $w$ an étale morphism, $\pi$ a finite surjective morphism of finite presentation, and $\phi^{-1}(s)$ consists of a single point $s''$, such that $h(s'')=x$ and $\kappa(s'')=\kappa(x)$.\nde{Consequently, $h$ induces a section $\sigma$ of $X\times_SS''\to S''$ such that $\sigma(s'')$ lies above $x$; compare with EGA IV$_3$, 14.5.10.}
\end{lemma}
\begin{proof}
First one may suppose $S=\Spec A$ and $X=\Spec B$, where $B$ is an $A$-algebra of finite presentation. Let $\mathfrak p$ and $\mathfrak q$ be the prime ideals of $A$ and $B$ corresponding to $s$ and $x$, respectively, so that $\mathfrak p=A\cap\mathfrak q$. Put $n=\dim_xX_s$, and let $t_1,\ldots,t_n$ be elements of $B$ whose images in $\OO_{X_s,x}=B_{\mathfrak q}/\mathfrak pB_{\mathfrak q}$ form a system of parameters. Then $\OO_{X_s,x}/(t_1,\ldots,t_n)$ is finite-dimensional over $\kappa(x)$ and hence also over $\kappa(s)$, since $x$ is a closed point of the algebraic $\kappa(s)$-scheme $X_s$. Consequently, the $S$-morphism
\[
u:X\longrightarrow\mathbb A_S^n=\Spec(A[T_1,\ldots,T_n]),
\]
defined by $T_i\mto t_i$, is of finite presentation, $x$ is isolated in its fiber $u^{-1}(u(x))$, and one has $u(x)=\tau_0(s)$, where $\tau_0:S\to\mathbb A_S^n$ denotes the ``zero section'' of $\mathbb A_S^n\to S$, corresponding to the morphism of $A$-algebras $A[T_1,\ldots,T_n]\to A$ which sends every $T_i$ to $0$. Since the set of points of $X$ isolated in their fiber above $\mathbb A_S^n$ is open (EGA IV$_3$, 13.1.4), one may suppose, shrinking $X$ if necessary, that $u$ is quasi-finite and $u^{-1}(u(x))=\{x\}$. By Lemma 5.6.1.0, $u$ is universally open at the point $x$.

Let $B_0=B/(t_1,\ldots,t_n)$ and let $X_0=X\times_{\mathbb A_S^n}\tau_0(S)=\Spec B_0$ and $u_0:X_0\to S$ be the morphism obtained from $u$ by the base change $\tau_0:S\hookrightarrow\mathbb A_S^n$. Then $u_0$ is quasi-finite and of finite presentation, universally open at the point $x$, and $x$ is the unique point of $X_0$ above $s$.

Let $A'$ be the henselization of the local ring $A_{\mathfrak p}=\OO_{S,s}$, and let $S'=\Spec(A')$, $B'_0=B_0\otimes_AA'$, and $X'_0=X_0\times_SS'=\Spec(B'_0)$. Then the closed point $s'$ of $S'$ is the unique point of $S'$ above $s$, one has $\kappa(s')=\kappa(s)$, and $X'_0$ has a unique point $x'$ above $s'$; one has $\kappa(x')=\kappa(x)$ and $x'$ is also the unique point of $X'_0$ above $s$ and $x$. Since $A'$ is henselian, by EGA IV$_4$, 18.5.11, $X'_0$ is the disjoint sum of two open and closed subsets:
\[
\tag{*}
X'_0=V\sqcup W,\qquad\text{where }V=\Spec(\OO_{X'_0,x'});
\]
and the local ring $\OO_{X'_0,x'}$ is finite and of finite presentation over $A'$. The restriction $\pi$ of $u'_0$ to $V$ is therefore finite and of finite presentation. Moreover, since $u'_0:X'_0\to S'$ is open at $x'$, one has $u'_0(V)=\Spec(\OO_{S',s})=S'$, so $\pi$ is surjective.
% typo?: kept as printed: O_{S',s} rather than O_{S',s'}.

This proves the desired result when $S=S'$. In the general case, $A'$ is a filtering inductive limit of subalgebras $A_i$ étale over $A$, and such that $S_i=\Spec(A_i)$ has a unique point $s_i$ above $s$ (and one has $\kappa(s_i)=\kappa(s)$). Then $B'_0=\varinjlim B_i$, where $B_i=B_0\otimes_AA_i$. Put $X_i=\Spec(B_i)=X_0\times_SS_i$ and $C=\OO_{X'_0,x'}$. By (*) above, one has $C\simeq B'_0/fB'_0$, for some idempotent $f\in B'_0$, and there is an index $i$ and $f_i\in B_i$ such that $f$ is the image of $f_i$ in $B'_0$. Put $C_i=B_i/f_iB_i$ and $V_i=\Spec(C_i)$.

Then $C=C_i\otimes_{A_i}A'$, whence $V=V_i\times_{S_i}S'$, and $C_i$ is an $A_i$-algebra of finite presentation (since this holds for $B_i$). Consequently, the morphisms
\[
\xymatrix{X'_0\ar[dr]^{u'_0}&\\V\ar[u]^{h'}\ar[r]^\pi&S'}
\]
come, by the base change $S'\to S_i$, from morphisms of finite presentation:
\[
\xymatrix{X_i\ar[dr]^{u_i}&\\V_i\ar[u]^{h_i}\ar[r]^{\pi_i}&S_i}\ .
\]
For every $j\ge i$, let $V_j=V_i\times_{S_i}S_j$ and let $\pi_j:V_j\to S_j$ be the morphism (of finite presentation) obtained from $\pi_i$ by base change. Since $\pi:V\to S'$ is finite and surjective, by EGA IV$_2$, 8.10.5, there is an index $j$ such that $\pi_j:V_j\to S_j$ is finite and surjective. Then $w:S_j\to S$ is affine étale, $S_j$ has a unique point $s_j$ above $s$, and $V_j$ has a unique point $x_j$ above $s_j$ (since $x'$ is the unique point of $V=V_j\times_{S_j}S'$ above $s'$):
\[
\xymatrix@C=3em{X_j\ar[r]\ar[dr]^{u_j}&X_0\ar[r]\ar[dr]^{u_0}&X\ar[d]^f\ar[dr]^u&\\V_j\ar[u]^{h_j}\ar[r]^{\pi_j}&S_j\ar[r]^w&S\ar[r]^{\tau_0}&\mathbb A_S^n}\ .
\]
Thus $x_j$ is the unique point of $V_j$ above $s$, its image under the morphism $V_j\to X_j\to X$ is $x$, and one has $\kappa(x_j)=\kappa(x)$.
\end{proof}
\begin{lemma}{5.6.2.0}\label{VIB.5.6.2.0}\nde{This lemma, communicated by O. Gabber, has been added. It makes it possible to simplify the proof of 5.6.2, and to prove Theorem 5.3, as well as 5.4, in a more general form; see 5.7 and 5.8 below.}
Let $k$ be a field and $G$ a nonempty $k$-scheme of finite type.

(i) Let $K$ be an extension of $k$ and $W$ a dense open subset of $G_K$. Denote by $\pi$ the projection $G_K\to G$; then
\[
U=\{g\in G\mid W\text{ contains every maximal point of }\pi^{-1}(g)\}
\]
is a dense open subset of $G$. (N.B. If $g$ is a closed point of $G$ belonging to $U$, one therefore has $\pi^{-1}(g)\subset W$.)

(ii) Suppose moreover $G$ geometrically irreducible. Let $\mu:G\times X\to Y$ be a morphism of $k$-schemes, $x$ a point of $X$, and $\Omega$ an open subset of $Y$ such that $\mu_x^{-1}(\Omega_{\kappa(x)})\ne\varnothing$, where $\mu_x$ denotes the morphism $G_{\kappa(x)}\to Y_{\kappa(x)}$, $g\mto\mu(g,x)$. Then
\[
U=\{g\in G\mid\mu\text{ sends every maximal point of }g\times x\text{ into }\Omega\}
\]
is a dense open subset of $G$, and, for every closed point $g$ of $G$ belonging to $U$, one has $\mu(g\times x)\subset\Omega$ and hence $\mu(g',x)\in\Omega_{\kappa(x)}$ (resp. $\mu(g,x')\in\Omega_{\kappa(g)}$), for every point $g'\in G_{\kappa(x)}$ above $g$ (resp. $x'\in X_{\kappa(g)}$ above $x$).

(iii) Suppose that $G=H^0$, where $H$ is a $k$-group scheme locally of finite type, acting on a nonempty $k$-scheme $X$, so that the morphism $H\times_SX\to X\times_SX$, $(h,x)\mto(hx,x)$ is surjective. Let $U$ be an open subset of $X$. Then:
\begin{enumeratea}
\item $G\cdot U$ is an open subset of $X$, equal to the union of the irreducible components of $X$ whose generic point belongs to $U$.
\item Suppose $U$ dense in every irreducible component of $X$. Then, for every finite subset $A$ of $X$, there is a closed point $g\in G$ such that $g\cdot A'\subset U_{\kappa(g)}$, where $A'$ is the inverse image of $A$ in $X_{\kappa(g)}$. In particular, the morphism $G\times U\to X$ is surjective.
\end{enumeratea}
\end{lemma}
% typo?: kept as printed: products over S in (iii), although the base is k.
\begin{proof}
(i) Let $\overline k$ be a separable closure of $k$ and let $L$ be an extension of $k$ containing a copy of $\overline k$ and of $K$. Denote by $\pi_L$ (resp. $\phi$) the projection $G_L\to G$ (resp. $G_L\to G_K$). Since, for every $g\in G$, $\phi$ sends the maximal points of $\pi_L^{-1}(g)$ surjectively onto those of $\pi^{-1}(g)$, one sees that $U=\{g\in G\mid W_L\text{ contains every maximal point of }\pi_L^{-1}(g)\}$. Thus, replacing $K$ by $L$, one is reduced to the case where $K$ contains $\overline k$.

Since the projection $p:G_K\to G_{\overline k}$ is surjective and open, $V=p(W)$ is a dense open subset of $G_{\overline k}$, and, since $p^{-1}(g)=\Spec(\kappa(g)\otimes_{\overline k}K)$ is irreducible for every $g\in G_{\overline k}$ (cf. EGA IV$_2$, 4.3.3), for every $g\in V$ the generic point of $p^{-1}(g)$ belongs to $W$. On the other hand, let $\mathscr G=\Gal(\overline k/k)$; since the projection $q:G_{\overline k}\to G$ is surjective and $\mathscr G$ acts transitively on the fibers of $q$, $U=q(V')$, where $V'$ is the intersection of the $\mathscr G$-conjugates of $V$.

Now let $Z$ be the closed subset $G_{\overline k}-V$, endowed with the reduced closed subscheme structure. Since $G_{\overline k}$, and hence $Z$, is of finite presentation over $\overline k$, $Z$ comes by base change from a reduced closed subscheme $Z_1$ of $G\otimes_kk_1$, for some finite Galois extension $k_1$ of $k$, so the $\mathscr G$-conjugates of $Z$ are finite in number, so their union is still a nowhere dense closed subset $F$ of $G_{\overline k}$, and $V'=G_{\overline k}-F$ is a dense open subset of $G_{\overline k}$. Thus, since the projection $q:G_{\overline k}\to G$ is surjective and open, $U=q(V')$ is a dense open subset of $G$. Moreover, for every closed point $g$ of $G$, the fiber $\pi^{-1}(g)$ consists of a finite number of closed points of $G_K$, so if $g\in U$, then $\pi^{-1}(g)\subset W$. This proves (i), and (ii) follows from it.

On the other hand, (iii)(a) has been proved in VIA, 2.6.4. Finally, if $U$ is dense in every irreducible component of $X$, $X=G\cdot U$, so, for every $x\in X$, $\mu_x^{-1}(U_{\kappa(x)})$ is a nonempty open subset of $G_{\kappa(x)}$, and then (iii)(b) follows from (ii).
\end{proof}
\begin{corollary}{5.6.2}\label{VIB.5.6.2}
Let $S,G,X$ be as in the preliminary hypotheses of 5.3, and let $U$ be an open subset of $X$, $s\in S$, and $A$ a finite subset of $X_s$. Suppose $U_s$ dense in $X_s$ and $X_s$ locally of finite type over $\kappa(s)$.\nde{The hypothesis on $X_s$ has been added and the proof simplified (and corrected) by taking account of the addition 5.6.2.0. On the other hand, the proof shows that the conclusion holds if one supposes only that $U_s$ is dense in every irreducible component of $X_s$.}
% Corollary 5.6.2 continues in en-06.tex.
